\documentclass[11pt,twocolumn]{article}

\usepackage{array}
\usepackage[url=false, doi=false]{biblatex}
\usepackage{booktabs}
\usepackage[font=footnotesize, labelfont=bf]{caption}
\usepackage{enumitem}
\usepackage[symbol]{footmisc}
\usepackage[scale=0.85]{geometry}
\usepackage{lipsum}
\usepackage{microtype}
\usepackage{multirow}
\usepackage{parskip}
\usepackage{siunitx}
\usepackage{titlesec}
\usepackage[dvipsnames]{xcolor}
\usepackage{xurl}
\usepackage[hidelinks]{hyperref}
\usepackage[nameinlink,capitalise]{cleveref}

\addbibresource{../bibliography.bib}

\titlelabel{}

\newcommand{\atom}[1]{\colorbox{Gray!30}{\ttfamily #1}}
\newcommand{\citen}[1]{~\parencite{#1}}
\newcommand{\citeay}[1]{\textcite{#1}}

\setcounter{footnote}{1}

\title{Modifying LeanCoP Options On the Fly \\[3pt] \large (aka ``the PAAR paper'')}
\date{}
\author{}

\begin{document}
\twocolumn[
	\centering
	\vspace{1em}
	{\LARGE Modifying LeanCoP Options On the Fly \par} \vspace{3pt}
	{\large (aka ``the PAAR paper'') \par}
	\vspace{2em}
]
% \maketitle

\section{Introduction and Research}

LeanCoP inspired a large family of connection provers that improve performance and completeness by making a small change to the control mechanism\citen{Otten2023leanCoP20}.

LeanCoP 2.0\citen{ottenLeanCoP20} added 3 control options to the original LeanCoP: \atom{cut} and \atom{scut}, to restrict backtracking in connection calculus, and \atom{comp(J)}, which makes iterative deepening complete by restarting the search and increasing the maximum height of this tree 1 at a time until it reaches maximum height \atom{J}, and restarting from height 1 and with cuts disabled after exhausting the final tree.

This work gathers ideas from other implementations of connection provers from different domains.

\citeay{letzModelEliminationConnection2001} show that, in tableau search, it's preferable to have subgoals with fewer solutions to detect dead ends earlier.
Reordering clauses to prioritise shorter clauses is a valid principle, because the number of solutions is roughly the product over the clause's subgoals.

% \citeay{zomboriRoleEntropyGuiding2021} introduces an entropy measurement to show how ``confident'' each policy should be: higher entropy suggests that the next step should be exploration, while lower entropy suggests that it should be exploitation\citen{zomboriRoleEntropyGuiding2021}.
% Adding entropy regularisation improves training performance, and can make simpler models have better results than more complex ones.

\citeay{zomboriRoleEntropyGuiding2021} show that the entropy of a learned policy's action distribution is a key factor in guiding connection provers: policies whose entropy is too low tend to overcommit to narrow search paths, while policies whose entropy is too high spread effort too thinly.
I adapt this insight at runtime by measuring the Shannon entropy of heuristic action scores at each proof step and use low entropy as a signal to add a \atom{cut} to the current connection action.

\citeay{liReinforcementLearningBased2024} frame the reset choice in SAT solvers as a multi-armed bandit problem.
They show that recent search productivity should influence reset behaviour.
Similar restart strategies for SAT solvers are also surveyet in \citeay{HoldenSATQSAT}.

\section{Novel Idea}

Triggering the three control atoms in LeanCoP, \atom{cut}, \atom{scut}, and \atom{comp(I)}, dynamically during the tree search rather than at the beginning can allow a prover to solve a tableau tree in less time and fewer steps.

A possible solution would be to use reinforcement learning, like RLCop\citen{Kaliszyk2018RLTP} or learnCoP [{\scriptsize citation coming soon!}].
However, inference in these systems is slow and a heuristic-based solution might solve problems more quickly on a subset of TPTP problems.

\section{Implementation}

I use the existing connection prover implementation as a base for this code\citen{connections_2023}.
This allows me to use the existing code for hard things like Prolog parsing, and also to use PyCoP as a control group.

My code implements 6 different strategies.
\begin{description}[style=nextline, itemsep=-3pt, leftmargin=16pt]
	\item[FirstAction \emph{\mdseries (control)}] Run the first legal action (PyCoP default).
	\item[LeanCoPCut \emph{\mdseries (control)}] Run \atom{cut}, \atom{scut}, and \atom{comp(7)}, like LeanCoP 2.0.
	\item[HeuristicSelection] Choose reductions if possible, otherwise choose the extension with the shortest clause.
	\item[EntropyCut] Run the first legal action, but add a cut when heuristic action scores have low entropy.
	\item[AdaptiveReset] Resets the tree if it hits 5 dead-ends in a row.
	\item[SmartCut] Combines the last three strategies.
\end{description}

% \footnotetext{Better name TBD.}

\section{Preliminary Results}

\newcommand{\MS}{{\color{BrickRed}\bfseries max-steps}}
\newcommand{\TO}{{\color{OrangeRed}\bfseries timeout}}
\newcommand{\best}{\color{Green}}

\sisetup{
    table-format = 11.0,
	table-text-alignment = right,
}

\setlength{\defaultaddspace}{4pt}

\begin{table*}[t]
	\centering
	\scriptsize
	\begin{tabular}{@{}>{\ttfamily}l @{\hspace{15pt}} S S @{\hspace{30pt}} S[table-format=16.0] S S[table-format=13.0] S@{}}
		\toprule
		{\rmfamily\bfseries Problem} & {\bfseries FirstAction} & {\bfseries LeanCopCut} & {\bfseries HeuristicSelection} & {\bfseries EntropyCut} & {\bfseries AdaptiveReset} & {\bfseries SmartCut} \\
		\midrule
		GEO179+3      & 11987 &       4744  &       9629  &       11979 & \best 2546  &        \MS \\
		GEO181+3      & 11929 &       4778  &       9812  &       11925 &       \MS   &        \TO \\
 		\addlinespace
		HWV026-3      & 10174 &       4212  &       \MS   &       10174 &       10174 &        \MS \\
 		\addlinespace
		KRS095+1      & 552   &       504   &       952   &       548   &       552   &       1014 \\
 		\addlinespace
		LCL186-1      & 245   &       477   &       297   & \best 245   &       368   &        \TO \\
		LCL644+1.001  & 8     &       8     & \best 2     &       8     &       8     &          2 \\
		LCL679+1.001  & 1     &       1     &       1     &       1     &       1     &          1 \\
 		\addlinespace
		MSC015-1.005  & 25401 &       29037 &       26257 &       28251 & \best 7556  &       8816 \\
 		\addlinespace
		NUM291+1      & 32860 &       19936 &       \MS   &       32860 &       32860 &        \MS \\
		NUM300+1      & 11135 &       7471  &       18308 &       11135 &       11135 &      18306 \\
		NUM615+1      & \MS   &       42778 &       \MS   &       \MS   &       \MS   &        \MS \\
		NUM634+3      & \TO   &       \TO   &    \best 2  &       \TO   &       \TO   &          2 \\
 		\addlinespace
		SET062+3      & 24    &       24    &       26    &       24    &       24    &         26 \\
		SET175+3      & \MS   &       12401 &       \MS   &       \MS   &       1062  & \best 1031 \\
		SET603+3      & \MS   &       9953  &       \MS   &       \MS   &       \TO   & \best  853 \\
		SET626+3      & \MS   &       35231 &       \MS   &       \MS   &       \TO   &        \TO \\
		SET886+1      & \MS   &       24582 &       \MS   &       \MS   &       \TO   &        \TO \\
		SET904+1      & 105   &       101   &       140   &       101   &       105   &        140 \\
		SEU121+2      & \MS   &       34536 &       \MS   &       \MS   &       \TO   &        \TO \\
 		\addlinespace
		SWV232+1      & 54    &       54    &       \TO   &       54    &       54    &        \TO \\
		SWV248-2      & \TO   &       \TO   &       7016  & \best 6018  &       \TO   &        \TO \\
		SWV264-2      & 9     &       9     &       11    &       9     &       9     &         11 \\
 		\addlinespace
		SYN061-1      & 46    &       46    &       46    &       46    &       46    &         46 \\
		SYN582-1      & 15843 &       11650 &       40624 &       15430 &       \TO   &        \TO \\
		SYN663-1      & 5651  &       \MS   & \best 5588  &       5626  &       \TO   &        \MS \\
		SYN709-1      & 16642 &       5238  &       11201 &       16065 &       \MS   &        \MS \\
		SYN726-1      & \MS   &       \MS   & \best 4468  &       \MS   &       \TO   &        \MS \\
		\bottomrule
	\end{tabular}
	\caption{Amount of steps required to solve representative TPTP problems. {\best Green} numbers show best results that are significantly better than the control groups. Only 27 of these problems finished within the time and step limit; larger experiments with a higher time limit will give us more information.}
	\label{big_table}
\end{table*}

A subset of 247 ``representative'' TPTP problems was used for evaluation\citen{chvalovsky2021characteristic}.
The results with any significance are found in \cref{big_table}.

The comparer runs all 6 policies in parallel, with a 30-second time limit and a \num{50000} step limit.
A larger experiment, with the entirety of TPTP and a longer time limit, will come in the future.

The results show that all three policies have marked improvements in certain problems.
It's likely that with further exploration we might get even better results than the control groups.

\renewcommand*{\bibfont}{\small}
\printbibliography{}

\end{document}
